\chapter{完整例题、接口注释与可执行测试}

本章把前面各章的“调用测试”落到可执行程序。测试文件不依赖 PDF 中偶然出现的全局变量，均从
\texttt{muban.cpp} 的 C++17 地基开始追加；每个接口在代码旁写明输入含义、返回值、初始化和多测清空
要求。运行某个测试文件时，所有断言通过才打印 \texttt{PASS}，断言失败会直接终止并给出行号。

\section{统一卡片接口注释}

每个板子接入题目时，先在 \texttt{solve()} 中完成以下四件事：

\begin{enumerate}
\item 读入题目数据，明确下标是 \(0\)-based 还是 \(1\)-based；
\item 调用初始化接口（例如 \texttt{build/init/clear}），清空节点池、时间戳、懒标记和多测答案；
\item 调用主接口，记录返回值或读取对象的查询接口；
\item 输出题目要求的答案，不把上一次测试的全局状态带入下一组。
\end{enumerate}

代码中的接口注释统一采用下面的格式；复制模板时只需要把题目变量映射到这些字段：

\begin{lstlisting}[language={C++}]
// 接口：当前例题所调用函数的完整签名、参数含义与下标范围。
// 前置：需要 build/init 的数据、值域限制和不变量。
// 返回：答案含义；无解、不连通或空结构如何表示。
// 清空：多测必须重置哪些数组、节点池、时间戳和懒标记。
\end{lstlisting}

若某个代码块没有自己的题面，使用测试文件中的调用测试；调用测试不是伪样例，而是直接对接口
执行操作并用断言检查答案。正式题例则给出输入输出，便于把调用方式改写进题目的
\texttt{solve()}。

\section{测试矩阵与正式题例}

\small
\setlength{\tabcolsep}{2pt}
\begin{longtable}[]{@{}p{0.16\textwidth}p{0.23\textwidth}p{0.22\textwidth}p{0.22\textwidth}@{}}
\toprule
领域 & 覆盖接口 & 测试文件（均在 \texttt{tests/}） & 正式题例 / 边界 \\
\midrule
\endhead
\bottomrule
\endlastfoot
基础、数学、数据结构 &
折半搜索、差分、二分答案、滑动窗口、单调结构、快速幂、扩展欧几里得、筛法、组合数恒等式、树状数组、线段树、01-Trie、回滚并查集、LIS &
\texttt{core.cpp} &
空集合、单点、组合数越界、整段更新、重复值、负更新、无可行解；可对照 P1177、P3368、P3372。\\
图论与树 &
BFS、多源 BFS、0-1 BFS、Dijkstra、拓扑排序、二分图染色、SCC、桥、Kuhn、Dinic、LCA &
\texttt{graph.cpp} &
重复源点、不可达点、零边、环、平行边、增广链、反向残量、根节点 LCA；可对照 P4779、P3387、P3388、P3376。\\
字符串 &
KMP、Z 函数、Manacher、Booth、AC、SAM、PAM、后缀数组与 LCP &
\texttt{strings.cpp} &
空串、重叠匹配、后缀模式、克隆状态、重复回文；可对照 P3375、P3796、P3804、P5496。\\
最小生成树与重构树 &
Kruskal、Prim、阈值连通块、瓶颈路 &
\texttt{mst.cpp} &
单点、重边、环、不连通、同权边、跨森林查询；可对照 P3366。\\
动态规划与几何 &
LIS、LCS、01/完全背包、区间/数位 DP、Nim、叉积、凸包、面积、线段相交 &
\texttt{dp\_geometry.cpp} &
重复值、容量为零、空区间、数字 4、重复点、共线点、端点相交；可对照 P1020、P1048、P1880、P2742。\\
进阶数学与数据结构 &
高精度整数、线性基、扩展 CRT、矩阵快速幂、高斯消元、主席树、李超树 &
\texttt{advanced.cpp} &
正负数除法与余数、超 128 位整数、负底数模幂、相关向量、非互质同余、无解方程、重复值、空直线集和交叉直线。\\
\end{longtable}
\normalsize

表中的测试文件均位于仓库根目录的 \texttt{tests/}；完整文件名和编译命令见仓库
\texttt{README.md}。

\section{基础与数学：完整小题例}

下面的题例与 \texttt{refactored\_core.cpp} 中的断言一一对应。把每行调用放入
\texttt{solve()} 即可得到正式题解；返回值含义和边界处理写在函数上方的接口注释中。

\subsection{子集和、区间加与二分答案}

\begin{lstlisting}
题例 A：给 n 个非负数和上界 limit，求不超过 limit 的最大子集和。
输入
4
0 0
1 6
3 8 3 5 6
4 9 8 1 4 4
输出
0
0
8
9

题例 B：把正数数组分成至多 k 段，最小化最大段和。
输入
2
4 2
7 2 5 10
4 3
1 2 3 4
输出
14
4
\end{lstlisting}

\subsection{树状数组、线段树与回滚并查集}

\begin{lstlisting}
题例：初始数组 [1,2,3,4]，执行区间加和区间和查询。
输入
4 4
1 2 3 4
Q 1 4
A 2 3 5
Q 1 4
Q 2 3
输出
10
20
15

边界调用：
Fenwick(1) 只更新 1；线段树整段懒标记后查询子区间；
rollback(snapshot) 后，合并前的集合大小和连通性必须恢复。
\end{lstlisting}

\subsection{高精度整数四则运算}

读入两个十进制整数 \texttt{a,b}（保证 \texttt{b!=0}），依次输出和、差、积、向零截断的商与余数：

\begin{lstlisting}
输入
123456789012345678901234567890 97
输出
123456789012345678901234567987
123456789012345678901234567793
11975308534197530853419753085330
1272750402189130710322005854
52
\end{lstlisting}

负数还要检查恒等式 \texttt{a=(a/b)*b+a\%b}；需要数学意义的非负余数时，再执行
\texttt{r=(r+b)\%b}（前提是 \texttt{b>0}）。

\section{图论与树：完整小题例}

\subsection{最短路、拓扑与最大流}

\begin{lstlisting}
题例：有向非负权图，从 1 号点求最短路。
输入
4 5
1 2 5
1 3 1
3 2 1
2 4 1
3 4 4
输出
0 2 1 3

题例：网络流的两条不相交路径。
输入
4 4
1 2 3
1 3 2
2 4 3
3 4 2
输出
5
\end{lstlisting}

不可达点输出 \texttt{INF}；拓扑排序若处理点数小于 \(n\) 必须报告有环；Dinic 每次分层后
重置 \texttt{it}，增广时同步修改反向边。

\subsection{SCC、桥、匹配与 LCA}

\begin{lstlisting}
调用测试：
SCC: 1->2, 2->3, 3->1, 3->4       -> 2 个分量
桥: 1-2 的平行边 + 2-3             -> 只有 2-3 是桥
Kuhn: 左 1->{1,2}, 左 2->{2}       -> 匹配数 2
LCA: 1-{2,3}, 3-{4,5}              -> lca(4,5)=3
\end{lstlisting}

\section{字符串、动态规划与几何：完整小题例}

\begin{lstlisting}
字符串：模式串 [a, ab, bab]，文本 abab
总出现次数应为 5；逐模式次数应为 [2, 2, 1]。

动态规划：物品重量 [2,3,4]、价值 [3,4,5]、容量 5
01 背包答案为 7；容量 0 时答案为 0。

几何：正方形 (0,0),(1,0),(1,1),(0,1)
凸包顶点数为 4，面积为 1；两条对角线相交，平行的两条单位线段不相交。
\end{lstlisting}

\section{一键验收}

在仓库根目录执行下列命令；每一行都必须输出 \texttt{PASS}，任一断言失败都应先修复模板
或测试，而不是删掉断言：

\begin{lstlisting}
g++ -std=c++17 -O2 -Wall tests/refactored_core.cpp
./a.out
g++ -std=c++17 -O2 -Wall tests/refactored_graph.cpp
./a.out
g++ -std=c++17 -O2 -Wall tests/refactored_strings.cpp
./a.out
g++ -std=c++17 -O2 -Wall tests/refactored_mst.cpp
./a.out
g++ -std=c++17 -O2 -Wall tests/refactored_dp_geometry.cpp
./a.out
g++ -std=c++17 -O2 -Wall tests/refactored_advanced.cpp
./a.out
\end{lstlisting}

这六个文件构成大版本当前的可执行冒烟集；章节中的边界清单仍然是扩展题目时的回归
用例，新增变体必须先加入对应测试，再把接口注释和例题补到卡片。
